\SourceLeaf{397}%
in der Mitte von \(AC,BD\) und \(EF\). Die Coordinaten dieser neun Punkte, Abscissen-Linie und Anfangspunkt ganz willkürlich gewählt, bezeichne ich mit \(a,a'\quad b,b'\quad c,c'\quad d,d'\) u.\,s.\,w. Da nun die drei Punkte \(A,B,E\) in einer geraden Linie liegen, so findet zwischen ihren Coordinaten folgende Bedingungsgleichung statt:
\[
a(e'-b')+b(a'-e')+e(b'-a')=0
\]
und eben so hat man, da \(ADF,BCF,DCE\) gerade Linien sind
\[
\begin{aligned}
a(f'-d')+d(a'-f')+f(d'-a')&=0\\
b(c'-f')+c(f'-b')+f(b'-c')&=0\\
c(e'-d')+d(c'-e')+e(d'-e')&=0
\end{aligned}
\]
Addirt man diese vier Gleichungen zusammen, so erhält man
\[
(a+c)(e'+f'-b'-d')+(b+d)(a'+c'-e'-f')+(e+f)(b'+d'-a'-c')=0
\]
oder da offenbar
\[
\begin{array}{lll}
\tfrac12(a+c)=g,&\tfrac12(b+d)=h,&\tfrac12(e+f)=i\\
\tfrac12(a'+c')=g',&\tfrac12(b'+d')=h',&\tfrac12(e'+f')=i'
\end{array}
\]
ist,
\[
g(i'-h')+h(g'-i')+i(h'-g')=0
\]
welches die Bedingungs-Gleichung ist, dass \(G,H,I\) in Einer geraden Linie liegen.
\par\bigskip\centerline{\rule{.25\textwidth}{.3pt}}
